% !TEX root = ../main.tex
\chapter{MpsBuilder and Executable Cache}
\label{ch:mps_builder}

The MpsBuilder is the middle stage of \Lucid's compile path. It
consumes a TraceIR (\chref{ch:tracer}) and produces a
\code{MPSGraphExecutable} that can be run on inputs of the
recorded shape. The executable cache holds compiled executables
keyed by the shape signature, so that repeated calls with the
same shape skip the compile step.

This chapter describes the MpsBuilder's data flow, the per-op
emitter dispatch, the executable cache's design, and the
performance characteristics of the compile and execute paths.

\section{The MpsBuilder Class}
\label{sec:mps_builder:class}

\code{MpsBuilder} is the Objective-C++ class in
\code{lucid/\_C/compile/MpsBuilder.\{h,mm\}} that translates a
\code{Trace} into an \code{MPSGraph}. The interface:

\begin{lstlisting}[style=lucidcpp,language=C++]
class MpsBuilder {
public:
  explicit MpsBuilder(const Trace& trace);

  // Build the MPSGraph and return the compiled executable.
  CompiledExecutable build();

  // Diagnostic: dump the constructed MPSGraph structure.
  void dump_graph(std::ostream& out) const;
};
\end{lstlisting}

The constructor takes a trace; \code{build()} walks the trace
node by node, invoking the per-op emitter for each, and assembles
the resulting \code{MPSGraphTensor} graph. The output is wrapped
in a \code{CompiledExecutable}
(\secref{sec:mps_builder:executable}).

\subsection{Per-node dispatch}
\label{ssec:mps_builder:per_node}

The build loop:

\begin{lstlisting}[style=lucidcpp,language=C++]
for (const TraceNode& node : trace.nodes) {
  // Look up the emitter by op name
  auto emitter = OpEmitterTable::lookup(node.op_name);
  if (!emitter) {
    // Fallback: defer this op to eager
    fallback_nodes.push_back(node.id);
    continue;
  }
  // Invoke emitter to build MPSGraphTensor for this node
  std::vector<MPSGraphTensor*> inputs;
  for (size_t id : node.input_ids) {
    inputs.push_back(mps_tensors[id]);
  }
  std::vector<MPSGraphTensor*> outputs = emitter(graph, inputs,
                                                 node.attrs);
  for (size_t i = 0; i < outputs.size(); ++i) {
    mps_tensors[node.output_id(i)] = outputs[i];
  }
}
\end{lstlisting}

The pattern is uniform across all ops; the per-op variation lives
inside the emitter functions (\chref{ch:op_emitters}).

\section{The Emitter Table}
\label{sec:mps_builder:emitter_table}

\code{OpEmitterTable} is a hash map from op name to emitter
function pointer:

\begin{lstlisting}[style=lucidcpp,language=C++]
class OpEmitterTable {
public:
  static EmitterFn lookup(const std::string& op_name);
  static void register_emitter(const std::string& op_name,
                               EmitterFn fn);
};
\end{lstlisting}

The table is populated at engine initialisation by
self-registering emitters in \code{lucid/\_C/compile/OpEmitters/}.
Each emitter file declares:

\begin{lstlisting}[style=lucidcpp,language=C++]
namespace {
  struct AddEmitterReg {
    AddEmitterReg() {
      OpEmitterTable::register_emitter("add", emit_add);
    }
  } _add_reg;
}
\end{lstlisting}

The pattern matches the op-registry self-registration
(\chref{ch:op_registry}~Section~\ref{ssec:op_registry:self_register}).

\subsection{Currently 183 emitters}
\label{ssec:mps_builder:emitter_count}

\Lucid\ 3.5 ships 183 emitters covering the common operations:
elementwise, reductions, matmul, conv, normalisations,
activations, indexing primitives, and the NN-specific helpers
(LayerNorm, MultiHeadAttention, SoftmaxCrossEntropy).

The ops without emitters (notably linalg, FFT, data-dependent
output ops) trigger the fallback to eager execution.

\section{Fallback Handling}
\label{sec:mps_builder:fallback}

When the builder encounters an op without a registered emitter,
two strategies are possible:

\begin{enumerate}
  \item \textbf{Partial fallback}: extract the unsupported op (and
    its dependencies) from the compiled region, run them eagerly,
    feed results back into the compiled graph.
  \item \textbf{Full fallback}: give up on compiling the whole
    region, run everything eagerly.
\end{enumerate}

\Lucid\ implements partial fallback for cleanly-bounded cases
(an unsupported op with no downstream supported ops fed back to
the compiled region) and full fallback for tangled cases. The
decision is per-trace and uses a simple connected-components
analysis on the trace DAG.

\subsection{Strict mode}
\label{ssec:mps_builder:strict_mode}

\code{lucid.compile(model, strict=True)} disables both fallback
modes: any unsupported op raises
\code{LucidError("op '<name>' not supported by compile")} with
the trace position. The user can then either implement the
emitter (a few dozen lines per op for most cases) or rewrite the
model to avoid the op.

\section{The CompiledExecutable}
\label{sec:mps_builder:executable}

\code{CompiledExecutable} wraps the \code{MPSGraphExecutable} that
\MPSGraph\ produces, plus the metadata needed to dispatch:

\begin{lstlisting}[style=lucidcpp,language=C++]
struct CompiledExecutable {
  MPSGraphExecutable* executable;     // compiled MPSGraph
  std::vector<size_t> input_ids;       // trace node ids for inputs
  std::vector<size_t> output_ids;      // trace node ids for outputs
  std::vector<TraceNode> fallback_nodes; // for partial-fallback
  ShapeSignature shape_key;            // for cache lookup
};

class CompiledExecutable {
public:
  std::vector<TensorImpl> run(const std::vector<TensorImpl>& inputs);
};
\end{lstlisting}

The \code{run()} method:
\begin{enumerate}
  \item Bind input tensors to the \MPSGraph\ executable's input
    slots.
  \item Execute via \code{MPSGraphExecutable::runWithMTLCommandQueue}.
  \item For partial fallback, splice the eager-executed results
    into the appropriate positions.
  \item Return the output tensors.
\end{enumerate}

\section{The Executable Cache}
\label{sec:mps_builder:cache}

\code{ExecutableCache} maps shape signatures to compiled
executables:

\begin{lstlisting}[style=lucidcpp,language=C++]
class ExecutableCache {
public:
  CompiledExecutable* get_or_compile(
      const Trace& trace_template,
      const std::vector<TensorImpl>& inputs);

  void clear();
  size_t size() const;
  size_t max_size() const;
  void set_max_size(size_t);
};
\end{lstlisting}

\subsection{The shape signature}
\label{ssec:mps_builder:shape_signature}

The cache key is a hash of the inputs' shapes and dtypes:

\begin{lstlisting}[style=lucidcpp,language=C++]
struct ShapeSignature {
  std::vector<std::pair<Shape, Dtype>> input_specs;

  bool operator==(const ShapeSignature& o) const;
  size_t hash() const;
};
\end{lstlisting}

Different shapes (different batch size, different sequence
length) produce different signatures and require separate
compilations.

\subsection{Cache size and eviction}
\label{ssec:mps_builder:cache_eviction}

The default maximum cache size is 32 entries. Beyond that,
least-recently-used (LRU) eviction removes the oldest entry.
The eviction is rare for typical training (the working set is
one or two shape signatures); it matters for inference workloads
that vary input shapes per call.

\subsection{Cache statistics}
\label{ssec:mps_builder:cache_stats}

\code{compiled\_module.cache\_stats()} returns hit/miss counts
and the current cache size. Useful for diagnosing whether a
training loop is hitting the cache (it should be) or
re-compiling per call (the user should investigate why shapes
are varying).

\section{The Compile Path Cost}
\label{sec:mps_builder:cost}

The cold compile of a non-trivial trace has three sequential
costs:

\begin{enumerate}
  \item Trace IR walking and MpsBuilder dispatch: tens of
    milliseconds for a typical training step.
  \item \MPSGraph\ graph construction: also tens of milliseconds.
  \item \MPSGraph\ compilation to executable: hundreds of
    milliseconds to seconds. This is the dominant cost.
\end{enumerate}

The third cost is opaque to us; \MPSGraph\ performs internal
optimisations (operator fusion, layout selection, kernel
selection) whose duration is workload-dependent.

\subsection{Cache hit cost}
\label{ssec:mps_builder:hit_cost}

A cache hit dispatches in roughly $50\,\mu s$: hash the input
shapes, look up in the cache map, return the executable. The
subsequent \code{executable->run()} adds another $\sim 100\,\mu s$
to set up the command buffer and submit.

\section{Per-Op MPSGraph Dispatch (3.4 Carve-out)}
\label{sec:mps_builder:per_op_carveout}

The 3.4 release introduced a per-op carve-out: certain ops, even
in eager mode, dispatch through MpsBuilder for a one-op compiled
graph that gives substantial speedup over the equivalent MLX
primitive. The current list includes GELU, LayerNorm, embedding
backward, max-pool backward, SiLU, and SoftmaxCrossEntropy
(\chref{ch:mpsgraph_dispatch}).

The infrastructure is identical to the full compile path: trace,
build, compile, cache. The difference is that the trace contains
exactly one op (the carve-out op), and the cache key includes
the eager-call's input shape.

\subsection{Why per-op}
\label{ssec:mps_builder:why_per_op}

Some operations have \MPSGraph\ implementations substantially
faster than \MLX's eager versions. The per-op carve-out lets
those ops benefit even in eager mode, without requiring the user
to call \code{compile()}. The result: eager training of a model
that uses GELU and LayerNorm (Transformer-typical) is roughly
20\% faster on the 3.4 release than on 3.3.

\section{Implementation Notes}
\label{sec:mps_builder:impl_notes}

The MpsBuilder is Objective-C++ (\code{.mm}) rather than pure
C++ because the \MPSGraph\ API is Objective-C. The bridge between
\Lucid's C++ engine and the Objective-C API is encapsulated in
the MpsBuilder file; the rest of the engine sees only the C++
\code{CompiledExecutable} wrapper.

\subsection{Memory management}
\label{ssec:mps_builder:mem_mgmt}

\code{MPSGraphExecutable} is a reference-counted Objective-C
object. \code{CompiledExecutable} holds a strong reference; the
\code{ExecutableCache} holds \code{CompiledExecutable}s; the
\code{CompiledModule} (\chref{ch:compiled_module}) holds the
cache. When the module is destroyed, the chain releases.

For very long-running processes that compile many distinct
graphs, the user can call \code{compiled\_module.empty\_cache()}
to release everything.

\section{Summary and Forward References}
\label{sec:mps_builder:summary}

The MpsBuilder consumes a TraceIR and produces a compiled
\MPSGraph\ executable, dispatched per-op through the
\code{OpEmitterTable}. The \code{ExecutableCache} keys by input
shape signature, supporting both the full \code{compile(model)}
path and the 3.4 per-op carve-out. Cold compile is hundreds of
milliseconds; cache hits are sub-millisecond.

The next chapter, \chref{ch:op_emitters}, catalogues the 183
per-op emitters and describes how a new emitter is added.
\chref{ch:fwd_bwd_compile} treats the forward+backward
unification. \chref{ch:compiled_module} treats the user-facing
\code{CompiledModule} class.

\begin{lucidnote}
For the user: the MpsBuilder is invisible at the API surface.
\code{compile()} returns a \code{CompiledModule}; the user calls
it like a normal module. The cache hit/miss is automatic; the
user only sees it if they call \code{cache\_stats()} for
diagnostics.
\end{lucidnote}
